% ─────────────────────────────────────────────────────────────────────────────
\section{Test Plan Summary and Results}
\label{sec:testing}
% ─────────────────────────────────────────────────────────────────────────────

\subsection{Reconstructed Validation Scope}

The retained validation cases covered five areas:

\begin{itemize}[nosep]
  \item input and detection,
  \item target isolation and semantic reasoning,
  \item output generation and local operation,
  \item deployment-oriented runtime behavior,
  \item known scope boundaries of the prototype.
\end{itemize}

\subsection{Test Cases}

\setlength{\LTleft}{0pt}
\setlength{\LTright}{0pt}
\begin{longtable}{@{}>{\ttfamily\small}p{1.5cm}p{5.4cm}p{7.2cm}}
  \caption{Final prototype test cases}\\
  \toprule
  ID & Test description & Expected result \\
  \midrule
  \endfirsthead

  \toprule
  ID & description & result \\
  \midrule
  \endhead

  TC-01 & Local thermal frame or stream ingestion & Frames are accepted from a local
  source and forwarded to the pipeline without external service dependence. \\
  TC-02 & Multi-object thermal detection & The detector identifies relevant targets such as
  people or vehicles and returns bounding boxes with coarse labels. \\
  TC-03 & Region-of-interest extraction & Each detection is cropped correctly and isolated
  for second-stage reasoning. \\
  TC-04 & Vehicle semantic analysis & The reasoner returns a richer interpretation than the
  plain label, including likely category and heat-related status clues. \\
  TC-05 & Human or armed-target semantic analysis & The reasoner produces a textual
  interpretation of posture or carried-object clues when present. \\
  TC-06 & Structured JSON generation & Detector outputs and reasoner outputs are merged
  into machine-readable structured records. \\
  TC-07 & Local presentation or logging & Final outputs are available for local review
  through logs, dashboard-style presentation, or both. \\
  TC-08 & Closed-network execution path & Runtime operation does not require external
  internet connectivity or third-party API calls. \\
  TC-09 & Containerized runtime startup & The packaged environment initializes with the
  required local dependencies and runtime services for prototype execution. \\
  TC-10 & Local model initialization & Detector and reasoning components load locally
  without remote service fallback before frame processing begins. \\
  TC-11 & No-detection handling and conditional execution & Frames with no relevant
  detection do not crash the pipeline and do not force unnecessary second-stage reasoning. \\
  \bottomrule
\end{longtable}

The deployment-related cases above They verify that the
system can be started, initialized, and operated locally as an edge-oriented workflow.

\subsection{Observed Bugs and Weaknesses}

The most significant weakness observed in the reconstructed validation record was the
inconsistency of very fine-grained recognition from thermal crops. This is not surprising: the
project deliberately operates in a domain with weak texture and limited identifying detail.

\subsection{Potential Enhancements}

The test results suggest clear future improvements:

\begin{itemize}[nosep]
  \item expand thermal-domain fine-tuning data for better object-specific reasoning,
  \item constrain prompts and output schemas to reduce over-general interpretation,
  \item preserve automated test scripts and benchmark logs as first-class project artifacts,
  \item formalize Jetson-oriented profiling for FPS, RAM usage, and startup stability,
  \item perform longer-duration embedded and field-style evaluation on the final target
    hardware.
\end{itemize}
